\chapter{Tangent Assembly}

This is the chapter where the PETSc redesign differs most strongly from the old
MATLAB implementation. The mathematical operator is still
\[
K_{\mathrm{tangent}} = B^T D_S B,
\]
but the default 3D PMG path no longer materializes that object as one globally
rebuilt sparse matrix.

\section{Mainline Path In One Line}

For the committed benchmark, the normal tangent route is:
\[
\begin{aligned}
\text{owned elastic rows}
&\rightarrow \text{owned tangent pattern} \\
&\rightarrow \text{overlap constitutive data} \\
&\rightarrow \text{rows kernel} \\
&\rightarrow \text{regularized free operator for PMG shell}.
\end{aligned}
\]

That is the single path MATLAB authors should hold in mind first.

\section{Why This Path Is Selected}

\petscmod{cli/assembly\_policy.py} selects the owned distributed path because:
\begin{itemize}
\item \code{mpi\_distribute\_by\_nodes = true},
\item the solver is not \code{DIRECT},
\item \code{constitutive\_mode = overlap}.
\end{itemize}

As a result:
\begin{itemize}
\item the runner enables the owned tangent path,
\item the runner enables the lightweight owned elastic path,
\item the global MATLAB-style tangent rebuild is not used in the normal solve.
\end{itemize}

\section{Step 1: Owned Elastic Rows}

The current path begins with
\code{assemble\_owned\_elastic\_rows\_for\_comm(...)} in
\petscmod{fem/distributed\_elastic.py}. This builds a PETSc-ready elastic
precursor by rank-owned rows.

\paragraph{What it stores.}
The owned elastic object stores:
\begin{itemize}
\item the node interval owned by the rank,
\item the overlap elements needed to assemble those rows,
\item a local CSR matrix containing only the rank-owned global rows,
\item the matching local right-hand-side contribution.
\end{itemize}

\paragraph{Why it matters.}
On the mainline path, the elastic matrix is not just a one-time startup object.
It becomes the structural template for later tangent regularization and for the
owned-row sparsity pattern.

\section{Step 2: Fixed Owned Tangent Pattern}

The core PETSc object is \code{OwnedTangentPattern} in
\petscmod{fem/distributed\_tangent.py}. For the current benchmark it is
prepared with:
\begin{itemize}
\item \code{include\_unique = false},
\item \code{include\_legacy\_scatter = false},
\item \code{include\_overlap\_B = false}.
\end{itemize}

Those three booleans are the clearest summary of what the current path is
\emph{not} doing.

\paragraph{What the pattern precomputes.}
For the mainline run, the pattern stores:
\begin{itemize}
\item the owned row interval and owned free-row mask,
\item the overlap node and element patch needed for those rows,
\item the fixed CSR sparsity pattern for those rows,
\item elastic values projected onto that same pattern,
\item overlap strain geometry data needed to re-evaluate the tangent each
Newton step,
\item constrained diagonal positions for Dirichlet rows.
\end{itemize}

This is the central architectural difference from MATLAB: the sparse pattern is
prepared once, and only the values are refreshed during Newton.

\section{Step 3: Overlap Constitutive Evaluation}

Because the mainline constitutive mode is \code{overlap}, each rank evaluates
stress and tangent data directly on the overlap integration points needed by its
owned rows. In practical terms:
\begin{itemize}
\item there is no unique-owner recovery stage,
\item there is no allgather of constitutive data,
\item the tangent kernel receives local overlap data directly.
\end{itemize}

That is why the overlap mode and the owned-row tangent pattern fit together so
well on the PMG path.

\section{Step 4: Row-Slot Tangent Kernel}

The mainline tangent kernel is \code{rows}. This is the normal production path
for the committed benchmark because \code{tangent\_kernel = rows} by execution
default and the benchmark does not override it.

\paragraph{How the rows kernel thinks.}
Instead of building a whole elemental matrix and scattering it afterward, the
kernel walks the fixed owned CSR pattern row by row. Its metadata is organized
around:
\begin{itemize}
\item \code{row\_slot\_ptr},
\item \code{slot\_elem},
\item \code{slot\_lrow},
\item \code{slot\_pos}.
\end{itemize}

So the assembly order is no longer ``form \(k_e\), then scatter.'' It is
``visit this owned sparse row, accumulate exactly the entries needed for this
row.'' The bilinear form is unchanged. The dataflow is row-oriented instead of
element-oriented.

\section{Step 5: Internal Force And Tangent Refresh}

At each Newton iteration the mainline path performs two distributed refreshes:
\begin{itemize}
\item \code{assemble\_owned\_force\_from\_local\_stress(...)} updates the owned
internal-force rows,
\item \code{assemble\_owned\_tangent\_values(...)} refreshes the tangent values
on the fixed owned CSR pattern.
\end{itemize}

This is the PETSc replacement for rebuilding the global
\code{B' * D\_p * B} matrix in MATLAB. The matrix structure stays fixed. Only
the values derived from the current constitutive state change.

\section{Step 6: Regularization On The Same Pattern}

The regularized Newton matrix is still
\[
K_r = r K_{\mathrm{elast}} + (1-r) K_{\mathrm{tangent}}.
\]

On the mainline path, that blend happens directly on the owned pattern through
\code{assemble\_owned\_regularized\_matrix(...)}:
\begin{itemize}
\item the elastic contribution comes from stored \code{elastic\_values},
\item the tangent contribution comes from the current owned tangent refresh,
\item the result keeps the same owned CSR structure.
\end{itemize}

This is why the elastic precursor is carried forward explicitly into the
nonlinear solve.

\section{Constrained Rows Stay Explicit}

Dirichlet rows are enforced through stored
\code{constrained\_diag\_positions}. After each tangent or regularized-matrix
refresh, those positions are reset to \(1\). So constrained rows remain valid
even though the unconstrained entries are overwritten every Newton step.

\section{What The PMG Shell Solver Receives}

The current PMG benchmark does not consume the whole operator family exposed by
the constitutive layer. The practical flow is:
\begin{enumerate}
\item owned elastic rows and owned tangent rows are built in PETSc format,
\item the regularized operator is formed on that owned-row pattern,
\item the free-space operator is extracted for the Krylov solve,
\item the PMG shell backend uses that reduced operator as its fine-grid system.
\end{enumerate}

So the default path is a distributed owned-row assembly underneath a reduced
free-space Newton solve.

\section{Where To Edit}

\begin{itemize}
\item \petscmod{cli/assembly\_policy.py} for path selection,
\item \petscmod{fem/distributed\_elastic.py} for the elastic precursor,
\item \petscmod{fem/distributed\_tangent.py} for the owned pattern and the
\code{rows} kernel,
\item \petscmod{constitutive/problem.py} for how the cached PETSc matrices are
exposed to Newton.
\end{itemize}

\section{Where The Other Tangent Paths Went}

Legacy scatter assembly, the global reference rebuild, unique-owner
constitutive recovery, and BDDC local-subdomain matrices are moved to
Appendix~\ref{app:tangent}.
